\documentclass[12pt,a4paper]{article}

\usepackage[margin=1.4cm]{geometry}
\usepackage{array}

\usepackage{enumitem}

% Global compact spacing for lists
\setlist[itemize]{
	topsep=0pt,
	partopsep=0pt,
	parsep=0pt,
	itemsep=0pt
}

\setlist[enumerate]{
	topsep=0pt,
	partopsep=0pt,
	parsep=0pt,
	itemsep=0pt
}

\usepackage{multicol}

\newcolumntype{L}[1]{>{\raggedright\arraybackslash}p{#1}}

\newcommand{\checkbox}{\fbox{\rule{0pt}{0.8ex}\rule{0.8ex}{0pt}}}
\usepackage{xcolor}
\definecolor{solutionred}{HTML}{7F1D1D}
\newcommand{\sol}[1]{{\color{solutionred}\small\strut #1\strut}}
\newcommand{\answerline}[1]{\makebox[0pt][l]{\raisebox{0.6ex}{\sol{#1}}}}
\begin{document}
	
	\begin{center}
		{\Large\bfseries C2: OBSERVAR EL CONTEXTO}
		
		\vspace{0.1cm}
		
		Emprendimiento e Innovación Grado 10\\[-1mm]{\scriptsize\color{solutionred}Modelo pedagógico: toda evidencia y respuesta es simulada; no son datos del colegio.}
		
		\vspace{0.3cm}
		
		\begin{tabular}{L{0.13\textwidth}L{0.65\textwidth}}
			Grado: & \answerline{10}\rule{2cm}{0.4pt} \\
			Fecha: & \rule{2cm}{0.4pt} \\
			Nombre 1: & \rule{\linewidth}{0.4pt} \\
			Nombre 2: & \rule{\linewidth}{0.4pt} \\
			Nombre 3: & \rule{\linewidth}{0.4pt} \\
		\end{tabular}
	\end{center}
	
	\section*{Nuestro objetivo}
	
	Ya han identificado varias oportunidades de mejora. Ahora, su equipo seleccionará una oportunidad e investigará qué está sucediendo realmente en el contexto donde ocurre.
	
	El objetivo no es confirmar que su idea original es correcta. El objetivo es comprender mejor la situación mediante una observación directa y sistemática.
	
	En esta etapa, presten atención a lo que hacen las personas, a lo que sucede a su alrededor, a cómo el entorno afecta la situación y a si ciertos patrones aparecen repetidamente.
	
	\section*{Seleccionen una oportunidad}
	
	Revisen las oportunidades identificadas en la Clase 1. Seleccionen una situación que su equipo considere que vale la pena investigar más a fondo.
	
	\vspace{-1.5mm}\begin{center}
		\small
		\begin{tabular}{|L{0.28\textwidth}|L{0.55\textwidth}|}
			\hline
			\textbf{Elemento} & \textbf{Nuestra selección} \\
			\hline
			\textbf{Oportunidad} & \rule{0pt}{0.6cm}\begin{minipage}[b]{\linewidth}\sol{Filas para comprar alimentos, seleccionadas en C1.}\end{minipage} \\
			\hline
			\textbf{Lugar} & \rule{0pt}{0.6cm}\begin{minipage}[b]{\linewidth}\sol{Un punto de venta de la cafetería del colegio.}\end{minipage} \\
			\hline
			\textbf{Personas involucradas} & \rule{0pt}{0.6cm}\begin{minipage}[b]{\linewidth}\sol{Estudiantes compradores y personal de atención.}\end{minipage} \\
			\hline
			\textbf{¿Por qué observar esta situación?} & \rule{0pt}{0.8cm}\begin{minipage}[b]{\linewidth}\sol{Medir la espera y describir el servicio sin asumir que toda fila es innecesaria.}\end{minipage} \\
			\hline
		\end{tabular}
	\end{center}\vspace{-1.5mm}
	
	\vspace{0.2cm}
	
	\textbf{Importante:} No intenten demostrar que la oportunidad es real. Observen la situación de manera abierta y registren lo que realmente sucede.
	
	\section*{Qué observar}
	
	Al observar el contexto seleccionado, presten atención a diferentes aspectos de la situación:
	
	\begin{multicols}{2}
		\begin{itemize}[leftmargin=0.6cm,itemsep=0pt]
			\item Personas
			\item Acciones y comportamientos
			\item Interacciones
			\item Espacios físicos
			\item Objetos y materiales
			\item Uso de recursos
			\item Movimientos y rutinas
			\item Esperas o interrupciones
			\item Dificultades u obstáculos
			\item Situaciones repetidas
		\end{itemize}
	\end{multicols}
	
	Observen qué sucede antes, durante y después de la situación que están investigando. \\
	
	
	\textbf{Regla:} Registren primero las observaciones como evidencia. Las interpretaciones pueden agregarse por separado después de haber recopilado la evidencia.
	
	\section*{Método de observación}
	
	Antes de comenzar, decidan cómo observará su equipo la situación.
	
	\vspace{-1.5mm}\begin{center}
		\small
		\begin{tabular}{|L{0.30\textwidth}|L{0.55\textwidth}|}
			\hline
			\textbf{Elemento} & \textbf{Nuestro plan de observación} \\
			\hline
			\textbf{¿Dónde observaremos?} & \rule{0pt}{0.55cm}\begin{minipage}[b]{\linewidth}\sol{Fila, mostrador y salida del mismo punto de venta.}\end{minipage} \\
			\hline
			\textbf{¿Cuándo observaremos?} & \rule{0pt}{0.55cm}\begin{minipage}[b]{\linewidth}\sol{Antes, durante y después de un descanso; repetir en otros descansos.}\end{minipage} \\
			\hline
			\textbf{¿A qué prestaremos atención?} & \rule{0pt}{0.7cm}\begin{minipage}[b]{\linewidth}\sol{Llegadas, longitud de fila, elección, pago, entrega y cruces.}\end{minipage} \\
			\hline
			\textbf{¿Quién registrará las observaciones?} & \rule{0pt}{0.55cm}\begin{minipage}[b]{\linewidth}\sol{Un integrante cuenta; otro cronometra; otro describe acciones, sin nombres.}\end{minipage} \\
			\hline
			\textbf{¿Cómo las registraremos?} & \rule{0pt}{0.55cm}\begin{minipage}[b]{\linewidth}\sol{Notas anónimas; conteos cada 2 minutos y seguimiento de 3 compras.}\end{minipage} \\
			\hline
		\end{tabular}
	\end{center}\vspace{-1.5mm}
	
	\vspace{0.2cm}
	
	\textbf{Recuerden:} Siempre que sea posible, observen sin interferir en la situación. No inventen información que no puedan observar directamente.
	
	\section*{Qué deben registrar}
	
	Para cada observación, registren suficiente información para que otra persona pueda comprender lo que sucedió.
	
	\vspace{-1.5mm}\begin{center}
		\small
		\begin{tabular}{|L{0.17\textwidth}|L{0.24\textwidth}|L{0.46\textwidth}|}
			\hline
			\textbf{Elemento} & \textbf{Qué escribir} & \textbf{Ejemplo} \\
			\hline
			\textbf{Hora} & ¿Cuándo sucedió? 
			& 8:05 a. m., inicio de la clase de Matemáticas (10 de septiembre de 2026) \\
			\hline
			\textbf{Lugar} & ¿Dónde sucedió? 
			& Salón de grado 10, inmediatamente después del timbre \\
			\hline
			\textbf{Personas involucradas} & ¿Quiénes participaron? 
			& 28 estudiantes y 1 profesor \\
			\hline
			\textbf{Acción} & ¿Qué sucedió? 
			& Después del timbre, el profesor interrumpió la clase planeada para determinar qué estudiantes estaban presentes y registrar la asistencia. \\
			\hline
			\textbf{Evidencia} & ¿Qué vieron o escucharon directamente? 
			& El profesor dedicó aproximadamente 5 minutos a tomar la asistencia antes de continuar con la clase. En una clase de 50 minutos, esto representa el 10\% del tiempo disponible. \\
			\hline
			\textbf{Contexto} & ¿Qué más estaba sucediendo? 
			& Los estudiantes llegaron en momentos ligeramente diferentes. El profesor tuvo que revisar simultáneamente la lista de clase y a los estudiantes presentes en el salón. Un estudiante llegó después de que la asistencia ya había sido registrada y el profesor tuvo que actualizar el registro. \\
			\hline
			\textbf{Patrón / Pregunta} & ¿Qué podría valer la pena observar nuevamente? 
			& La asistencia se realiza durante la transición entre la entrada de los estudiantes al salón y el inicio de la clase. ¿Cuánto tiempo de clase se dedica a la asistencia y qué partes del proceso requieren mayor atención del profesor? \\
			\hline
		\end{tabular}
	\end{center}\vspace{-1.5mm}
	
	\vspace{0.2cm}
	
	\textbf{Importante:} Una observación útil es lo suficientemente específica como para que otra persona pueda comprender lo que sucedió sin haber estado presente.
	
	\section*{Observación}
	
	\vspace{-1.5mm}\begin{center}
		\small
		\begin{tabular}{|L{0.22\textwidth}|L{0.61\textwidth}|}
			\hline
			\textbf{Hora} & \rule{0pt}{0.45cm}\begin{minipage}[b]{\linewidth}\sol{Minutos 0--12 desde el inicio de un descanso simulado.}\end{minipage} \\
			\hline
			\textbf{Lugar} & \rule{0pt}{0.45cm}\begin{minipage}[b]{\linewidth}\sol{El mismo punto de venta de la cafetería seleccionado en C1.}\end{minipage} \\
			\hline
			\textbf{Personas involucradas} & \rule{0pt}{0.45cm}\begin{minipage}[b]{\linewidth}\sol{Estudiantes compradores; una persona atiende este punto en el ejemplo.}\end{minipage} \\
			\hline
			\textbf{¿Qué sucedió?} & \rule{0pt}{0.75cm}\begin{minipage}[b]{\linewidth}\sol{La fila aumenta y luego disminuye. Se pide, paga y recibe en el mostrador.}\end{minipage} \\
			\hline
			\textbf{Evidencia: ¿Qué vimos o escuchamos?} & \rule{0pt}{0.75cm}\begin{minipage}[b]{\linewidth}\sol{Fila en minutos 0/2/4/6/8/10/12: 4/10/16/18/15/11/7 personas. Tres seguimientos: 6, 7 y 5 minutos hasta iniciar atención; atención: 40, 55 y 45 segundos.}\end{minipage} \\
			\hline
			\textbf{Contexto} & \rule{0pt}{0.65cm}\begin{minipage}[b]{\linewidth}\sol{Dos compradores preguntan precios al llegar al mostrador; una entrega se detiene 30 segundos para reponer un producto. Dos personas cruzan junto a la fila; no se verifica pérdida de turnos.}\end{minipage} \\
			\hline
			\textbf{Interpretación / Pregunta} & \rule{0pt}{0.65cm}\begin{minipage}[b]{\linewidth}\sol{Hipótesis: llegadas concentradas y pausas podrían influir. Faltan tasas de llegada y atención para examinar qué espera es evitable.}\end{minipage} \\
			\hline
		\end{tabular}
	\end{center}\vspace{-1.5mm}
	
	\section*{Identifiquen patrones emergentes}
	
	Revisen sus observaciones. Busquen situaciones que hayan ocurrido más de una vez o condiciones que parezcan estar relacionadas con la oportunidad.
	
	\vspace{-1.5mm}\begin{center}
		\small
		\begin{tabular}{|m{0.2\textwidth}|m{0.7\textwidth}|}
			\hline
			\textbf{Pregunta} & \textbf{Nuestros hallazgos} \\
			\hline
			¿Qué ocurrió más de una vez? & \rule{0pt}{2.3cm}\begin{minipage}[b]{\linewidth}\sol{Dos consultas de precios al iniciar la atención y dos cruces junto a la fila. Son repeticiones dentro de una sola sesión simulada, no hábitos demostrados.}\end{minipage} \\
			\hline
			¿Cuándo ocurrió la situación con mayor frecuencia? & \rule{0pt}{2.3cm}\begin{minipage}[b]{\linewidth}\sol{La fila crece en los primeros seis minutos (4 a 18 personas) y después baja a 7. No tenemos evidencia para comparar días o descansos diferentes.}\end{minipage} \\
			\hline
			¿Quiénes fueron los más afectados? & \rule{0pt}{2.3cm}\begin{minipage}[b]{\linewidth}\sol{Los tres estudiantes seguidos esperaron 5 a 7 minutos antes de ser atendidos. El personal atiende una fila continua; su carga y la experiencia de otros usuarios deben consultarse.}\end{minipage} \\
			\hline
			¿Qué factores contextuales parecieron relevantes? & \rule{0pt}{2.5cm}\begin{minipage}[b]{\linewidth}\sol{Llegadas concentradas al comienzo, selección en el mostrador y reposición de un producto. Coinciden con la espera; no se ha medido la contribución causal de cada factor.}\end{minipage} \\
			\hline
			¿Qué nos sorprendió? & \rule{0pt}{2.3cm}\begin{minipage}[b]{\linewidth}\sol{La fila disminuye sin desaparecer y la atención sigue activa. Una fila larga no demuestra inactividad del personal ni que toda espera sea innecesaria.}\end{minipage} \\
			\hline
		\end{tabular}
	\end{center}\vspace{-1.5mm}
	
	
	\textbf{Un patrón no es simplemente una suposición.} Debe estar respaldado por observaciones recopiladas en el contexto.
	
	\newpage
	\section*{Reflexión inicial}
	
	\textbf{1. ¿Qué observaron que no esperaban encontrar?}
	
	\vspace{0.1cm}
	
	\noindent\answerline{La fila bajó de 18 a 7 personas entre los minutos 6 y 12, con atención activa.}\rule{\linewidth}{0.4pt}
	
	\vspace{0.35cm}
	
	\noindent\answerline{La espera puede variar dentro del mismo descanso; no equivale a falta de servicio.}\rule{\linewidth}{0.4pt}
	
	\vspace{0.4cm}
	
	\textbf{2. ¿Qué observación proporciona la evidencia más sólida sobre la situación? ¿Por qué?}
	
	\vspace{0.1cm}
	
	\noindent\answerline{Tres seguimientos separan espera (6, 7 y 5 min) de atención (40, 55 y 45 s).}\rule{\linewidth}{0.4pt}
	
	\vspace{0.35cm}
	
	\noindent\answerline{Son registros concretos, aunque insuficientes para estimar una espera típica.}\rule{\linewidth}{0.4pt}
	
	\vspace{0.4cm}
	
	\textbf{3. ¿Qué patrón o comportamiento repetido identificaron?}
	
	\vspace{0.1cm}
	
	\noindent\answerline{Dos compradores preguntaron precios al llegar al mostrador.}\rule{\linewidth}{0.4pt}
	
	\vspace{0.35cm}
	
	\noindent\answerline{Es una repetición de la sesión simulada; no prueba la causa principal de la fila.}\rule{\linewidth}{0.4pt}
	
	\vspace{0.4cm}
	
	\textbf{4. ¿Qué necesitan observar o comprender todavía?}
	
	\vspace{0.1cm}
	
	\noindent\answerline{Faltan otros descansos, tasas de llegada y atención y el tiempo restante para comer.}\rule{\linewidth}{0.4pt}
	
	\vspace{0.35cm}
	
	\noindent\answerline{Preguntaremos a estudiantes y personal qué dificultades viven y qué límites existen.}\rule{\linewidth}{0.4pt}
	
	\vspace{0.4cm}
	
	\textbf{5. ¿Algo de lo que observaron cuestionó sus suposiciones iniciales?}
	
	\vspace{0.1cm}
	
	\noindent\answerline{Sí: había atención continua y la fila disminuyó; no toda espera parece evitable.}\rule{\linewidth}{0.4pt}
	
	\vspace{0.35cm}
	
	\noindent\answerline{La reposición podría ser necesaria. Debemos comprenderla antes de juzgarla.}\rule{\linewidth}{0.4pt}
	
	\section*{Verificación final}
	
	\noindent\checkbox\makebox[0pt][r]{\sol{\(\times\)}\hspace{0.15em}}\quad Observamos un contexto real relacionado con la oportunidad seleccionada.
	
	\vspace{0.15cm}
	
	\checkbox\makebox[0pt][r]{\sol{\(\times\)}\hspace{0.15em}}\quad Registramos observaciones específicas en lugar de opiniones generales.
	
	\vspace{0.15cm}
	
	\checkbox\makebox[0pt][r]{\sol{\(\times\)}\hspace{0.15em}}\quad Identificamos a las personas, acciones y contexto involucrados.
	
	\vspace{0.15cm}
	
	\checkbox\makebox[0pt][r]{\sol{\(\times\)}\hspace{0.15em}}\quad Separamos las observaciones directas de las interpretaciones.
	
	\vspace{0.15cm}
	
	\checkbox\makebox[0pt][r]{\sol{\(\times\)}\hspace{0.15em}}\quad Buscamos situaciones repetidas o patrones emergentes.
	
	\vspace{0.15cm}
	
	\checkbox\makebox[0pt][r]{\sol{\(\times\)}\hspace{0.15em}}\quad Utilizamos evidencia para respaldar nuestras observaciones.
	
	\vspace{0.15cm}
	
	\checkbox\makebox[0pt][r]{\sol{\(\times\)}\hspace{0.15em}}\quad Identificamos lo que todavía necesitamos comprender.
	
	\vspace{0.3cm}
	
	\vspace{-1.5mm}\begin{center}
		\textbf{Primero observen lo que sucede. Luego interpreten lo que podría significar.}
	\end{center}\vspace{-1.5mm}

	
\end{document}
